% -------------------------------------------------------------
% SPRINTS
% -------------------------------------------------------------

\section{Sprint 1}
\textbf{Goal:} First sprint focusing on the beginning of the project, including report writing and initial coding tasks.

\medskip
\noindent \textbf{Short summary:} Both report and coding tasks were completed successfully. 2 bugs were not fixed. On report
entity framework still persisted. During this sprint the team realized the importance of correct story point assignment and the need for better communication between team members. The retrospective revealed that the team needs to improve on time management and task estimation for future sprints (Sprint 1 planning, backlog, and retrospective can be found in Appendix \ref{appendix:figures}).

\section{Sprint 2}

\medskip
\noindent \textbf{Short summary:} This sprint there were no report tasks. On the other hand the team made good progress on the code side,
getting closer to the first feasable demo. All tasks were successfully finished. Unit tests were written retroactively to Asset/Source services. The frontend was connected to the backend and Results/ResultLists were created. 
This sprint the team started experimenting with user stories that could not be finished in one sprint, because of significant complexity.
In cases like this subtasks were assigned as if they were normal tasks. These large items were left in TODO for each sprint until all
subtasks were done. Team members started figuring out the optimizing algorithm used by the system, and created the Asset Manager Window (Sprint 2 planning, backlog, and retrospective can be found in Appendix \ref{appendix:figures}).

\section{Sprint 3}
\textbf{Goal:} Third sprint focusing on the completion of the first demo, including the implementation of the optimization algorithm and the creation of a user interface.

\medskip
\noindent \textbf{Short summary:} The team managed to complete all required tasks for the demo, even though one task had to be left 
for the next sprint, as it was not necessary for the demo and it took up precious time from 
working on the demo. Minor tasks had to be assigned in the middle of the sprint. For example CSV creation for the demo was not planned, but it was necessary to show the results of optimization (Sprint 3 planning, backlog, and retrospective can be found in Appendix \ref{appendix:figures}).

\section{Sprint 4}
\textbf{Goal:} Fourth sprint focusing on LaTex report writing and implementing feedback received from the midterm evaluation.

\medskip
\noindent \textbf{Short summary:} The team managed to implement most of the feedback received on the midterm evaluation and the project
is wrapping up nicely. Members are going to focus on implementing extra 
requirements. Report got migrated from Word over to LaTeX so the team focused more on fitting it to the given template (Sprint 4 planning, backlog, and retrospective can be found in Appendix \ref{appendix:figures}).

\section{Sprint 5}
\textbf{Goal:} Fifth sprint focusing on report tasks as well as on small quality of life UI changes.

\medskip
\noindent \textbf{Short summary:} The first week of this sprint was unfortunate as not all team members were in Denmark and could not fully concentrate. Despite the short time, the team managed to implement all tasks preparing for final 
sprint. Except for merge conflicts slipping through and making the main branch not 
working, there were not any big problems (Sprint 5 backlog and retrospective can be found in Appendix \ref{appendix:figures}).

\section{Sprint 6}
\textbf{Goal:} Sixth sprint focusing on report and code finalization.

\medskip
\noindent \textbf{Short summary:} The team managed to finish all tasks for the final sprint, including report finalization and code cleanup. Overall, the project was a success and the team is proud of what they have accomplished (Sprint 6 planning and retrospective can be found in Appendix \ref{appendix:figures}).
